\section{R ile çözümlü uygulama}
\label{sec:r-b05}

\textbf{Soru.} Ortalaması ve standart sapması 10 olan üstel bir evrenden
$n=1,5,30$ hacimlerinde bağımsız örneklemler üretin. Örneklem ortalamalarının
merkezini ve standart hatasını karşılaştırın.

\begin{lstlisting}[style=prismR]
set.seed(2026)
tekrar <- 10000
hacimler <- c(1, 5, 30)
sonuc <- data.frame(n = hacimler,
                    ortalama = NA_real_,
                    ampirik_se = NA_real_,
                    kuramsal_se = 10 / sqrt(hacimler))
for (sira in seq_along(hacimler)) {
  hacim <- hacimler[sira]
  ortalamalar <- replicate(tekrar,
    mean(rexp(hacim, rate = 1 / 10)))
  sonuc$ortalama[sira] <- mean(ortalamalar)
  sonuc$ampirik_se[sira] <- sd(ortalamalar)
  hist(ortalamalar, breaks = 35, probability = TRUE,
       main = paste("n =", hacim), xlab = "Ortalama")
  abline(v = 10, col = "blue", lwd = 2)
}
print(sonuc)
c(se_25 = 10 / sqrt(25), se_100 = 10 / sqrt(100))
\end{lstlisting}

\begin{outputguidebox}
\textbf{Çözüm ve kuramsal kontrol.} Üç ortalama da yaklaşık 10 çevresinde
olmalıdır. Kuramsal standart hatalar sırasıyla 10, 4,4721 ve 1,8257'dir;
benzetimdeki standart sapmalar bu değerlere yaklaşır, tam eşit olmak zorunda
değildir. $n=25$ ve $n=100$ için standart hatalar 2 ve 1'dir.
\end{outputguidebox}

\textbf{Yorum.} R'deki üstel dağılımın \texttt{rate} parametresi ortalamanın
tersidir. \texttt{rate = 10} yazmak farklı bir evren üretir. Ham evren
çarpık kalırken ortalamaların dağılımı örneklem hacmi arttıkça daha simetrik
ve dar hale gelir. $n=30$ her evren için otomatik normallik güvencesi değildir.

\textbf{Pekiştirme ve yanıt.} Standart hatayı yarıya indirmek için hacim
dört katına çıkarılır. Aynı tohumun R ve Python'da aynı rastgele sayıları
üretmesi beklenmez; kuramsal hedefleri karşılaştırın.